\documentclass[11pt]{amsart}
\usepackage{amsmath,amssymb,stmaryrd}
\setlength{\emergencystretch}{3em}
\tolerance=600
\usepackage{tikz-cd}
\usepackage[hidelinks]{hyperref}

\newtheorem{thm}{Theorem}[section]
\newtheorem{lem}[thm]{Lemma}
\newtheorem{cor}[thm]{Corollary}
\newtheorem{prop}[thm]{Proposition}
\theoremstyle{definition}
\newtheorem{defi}[thm]{Definition}
\newtheorem{rem}[thm]{Remark}
\newtheorem{ex}[thm]{Example}
\newtheorem{que}[thm]{Question}

\newcommand{\Fix}{\mathrm{Fix}}
\newcommand{\Up}{\mathrm{Up}}
\newcommand{\Down}{\mathrm{Down}}
\newcommand{\Nuc}{\mathrm{N}}
\newcommand{\LLJ}{\mathbf{LJ}}
\newcommand{\PLL}{\mathbf{PLL}}
\newcommand{\PLLt}{\mathbf{PLL}_2}
\newcommand{\lax}{\bigcirc}
\newcommand{\laxa}{\bigcirc_1}
\newcommand{\laxb}{\bigcirc_2}
\newcommand{\Same}{\mathrm{Same}}
\newcommand{\Presh}{\mathrm{Presh}}

\begin{document}

\title[The lax logic of the join of two nuclei]{The join of two nuclei as a lax modality: axiomatization, finite sites, and inexpressibility}
\author{Ryota Shibusawa}
\thanks{Corresponding author. E-mail: \texttt{r-shibusawa@daiichi-koudai.ac.jp}}
\address{Daiichi Institute of Technology}
\email{r-shibusawa@daiichi-koudai.ac.jp}
\keywords{lax logic, nucleus, Lawvere--Tierney topology, Grothendieck topology,
  join of nuclei, intuitionistic modal logic, common knowledge, finite poset}
\subjclass[2020]{03B45 (primary), 03B20, 06D20, 18F10, 18B25 (secondary)}

\begin{abstract}
Propositional lax logic is the intuitionistic modal logic of a nucleus on a
Heyting algebra, equivalently of a Lawvere--Tierney topology on a topos. Two
nuclei have a join in the lattice of nuclei, whose fixed points are the common
fixed points; on a finite algebra it is a finite iterate of the composite, on a
site it is the join of the two Grothendieck topologies. We axiomatize the
trimodal logic of two lax modalities together with their join by adding to the
lax axioms the inclusions $\laxa\varphi\to\lax\varphi$, $\laxb\varphi\to\lax\varphi$
and one induction scheme,
$(\laxa\psi\to\psi)\to(\laxb\psi\to\psi)\to(\varphi\to\psi)\to(\lax\varphi\to\psi)$,
and prove algebraic completeness. On a finite poset $P$ the nuclei of $\Up(P)$
are the relativized necessities $\lax_C$ over subsets $C\subseteq P$, the join
of $\lax_C$ and $\lax_{C'}$ is $\lax_{C\cap C'}$, and the join equals the
$|P|$-fold iterate of the composite, which gives the finite-site axiom. The
join is nevertheless not expressible in the bimodal lax logic: on the chain
$\omega+2$ with the two nuclei stepping from even to odd and from odd to even
ordinals, every closed term of the two-nucleus language takes a value in
$\mathbb N\cup\{\omega+1\}$, while the join of the two nuclei sends $\bot$ to
$\omega$. The induction scheme is the lax analogue of the induction axiom of
common knowledge, with the join playing the role of the common-knowledge
operator and the iterates that of ``everybody knows'' to depth $n$.
\end{abstract}

\maketitle

\section{Introduction}

A \emph{nucleus} on a Heyting algebra $H$ is a map $j\colon H\to H$ with
$a\le ja$, $jja=ja$ and $j(a\wedge b)=ja\wedge jb$. Nuclei are the
Lawvere--Tierney topologies of topos theory restricted to the algebra of
subobjects of $1$, the ``geometric modalities'' of Goldblatt \cite{Goldblatt},
and the algebraic semantics of the propositional lax logic $\PLL$ of
Fairtlough and Mendler \cite{FM}: intuitionistic propositional logic with a
modality $\lax$ satisfying $\varphi\to\lax\varphi$,
$\lax\lax\varphi\to\lax\varphi$ and
$\lax\varphi\wedge\lax\psi\to\lax(\varphi\wedge\psi)$.

The nuclei of a complete Heyting algebra form a complete lattice under the
pointwise order; meets are pointwise, and the join $j_1\vee j_2$ is the nucleus
whose fixed points are the common fixed points of $j_1$ and $j_2$
\cite[II.2]{Johnstone}. The join is harder to compute: Macnab \cite{Macnab}
obtains it as a transfinite iterate of the composite $j_1j_2$, and on a finite
algebra it is a finite iterate. For Grothendieck topologies on a site the join
of nuclei is the join of topologies, so a statement about $j_1\vee j_2$ is a
statement about ``locally for $J\vee J'$''.

This paper treats the join as a modality. Let $\PLLt$ be the bimodal lax logic
with two lax modalities $\laxa,\laxb$, and let $\LLJ$ (Definition~\ref{def:LJ})
be the trimodal logic obtained by adding a third lax modality $\lax$ with the
axioms $\laxa\varphi\to\lax\varphi$, $\laxb\varphi\to\lax\varphi$ and the
\emph{induction scheme}
\[
\mathrm{(JL)}\qquad(\laxa\psi\to\psi)\to(\laxb\psi\to\psi)\to(\varphi\to\psi)\to(\lax\varphi\to\psi).
\]

\begin{thm}\label{thm:main}
\begin{enumerate}
\item[(a)] $\LLJ$ is algebraizable, its equivalent algebraic semantics being
the Heyting algebras with nuclei $j_1,j_2,j$ in which $ja$ is the least common
fixed point of $j_1$ and $j_2$ above $a$, for every $a$; on complete and on
finite Heyting algebras this says that $j$ is the join of $j_1$ and $j_2$ in
the lattice of nuclei, so every complete Heyting algebra with two nuclei and
their join, in particular every topos with two Lawvere--Tierney topologies and
their join, is such an algebra (Theorem~\ref{thm:complete}).
\item[(b)] On a finite poset $P$ the nuclei of $\Up(P)$ are the relativized
necessities $\lax_C\varphi$ $=$ ``$\varphi$ at every point of $C$ above the
current one'', $C\subseteq P$; the join of $\lax_C$ and $\lax_{C'}$ is
$\lax_{C\cap C'}$ and equals $(\lax_C\lax_{C'})^{|P|}$ (Theorem~\ref{thm:poset}).
Hence $\LLJ$ together with the scheme
$(\laxa\laxb)^{N+1}\varphi\to(\laxa\laxb)^N\varphi$ proves
$\lax\varphi\leftrightarrow(\laxa\laxb)^N\varphi$, and the scheme holds on
every poset with at most $N$ points.
\item[(c)] The join is not expressible in $\PLLt$: there is no formula $\chi$ in
$\laxa,\laxb$ such that $\lax\bot\leftrightarrow\chi$ holds in all
$\LLJ$-algebras, and a fortiori no $\chi(p)$ with
$\lax p\leftrightarrow\chi(p)$ (Theorem~\ref{thm:inexp}).
\end{enumerate}
\end{thm}

The induction scheme is the intuitionistic, inflationary counterpart of the
induction axiom of the logic of common knowledge \cite{HM}: common knowledge
$C\varphi$ is the greatest fixed point of $X\mapsto E(\varphi\wedge X)$, where
$E$ is ``everybody knows'', and is not a finite iterate $E^n\varphi$; the join
$\lax\varphi$ is the least fixed point of $X\mapsto\varphi\vee\laxa X\vee\laxb X$
(Lemma~\ref{lem:lfp}) and is not a finite iterate $(\laxa\laxb)^n\varphi$.

\subsection*{What is known.}
Lax logic and its algebraic, Kripke and cover semantics are in \cite{FM,
Goldblatt,GoldblattCover}. That the join of nuclei on a frame has as fixed
points the common fixed points is classical \cite[II.2]{Johnstone}; Macnab
\cite{Macnab} computes the join by transfinite iteration of the composite, and
Escard\'o \cite{Escardo} computes joins of arbitrary families of nuclei as
least fixed points of inflationary operators on prenuclei and deduces, from
Pataraia's fixed-point theorem, an induction principle for joins of nuclei.
The scheme (JL) is a formula-level internalization, inside the Heyting
algebra, of the common-fixed-point description of the join; what is added here
is the deductive system built on it and the study of its expressive power. Nuclei of the form
$\lax_C$ on the up-sets of a finite poset are the finite case of the
representation of nuclear Heyting algebras and nuclear implicative semilattices
by posets with a distinguished subset \cite{BH,BBCGGJ}. For a finite category the Grothendieck
topologies correspond to idempotent two-sided ideals (Kelly--Lawvere
\cite{KL89}, Marqu\`es \cite{Marques25}), and the join topology to the largest
idempotent ideal in the intersection, reached by $2m$ alternating products
\cite{Finitejoin}; part (b) is the poset case of this dictionary.

\subsection*{What is new.}
The deductive system $\LLJ$, its algebraic completeness, the alternation bound
of Theorem~\ref{thm:poset}(ii) with its sharpness, and the inexpressibility
Theorem~\ref{thm:inexp}: the join of two geometric modalities is a genuinely
new connective, uniquely determined by the two but not definable from them,
in the same way as common knowledge is new relative to the individual knowledge
operators (Corollary~\ref{cor:beth}).

\section{Nuclei and their joins}\label{sec:nuclei}

Let $H$ be a Heyting algebra and $\Nuc(H)$ the set of its nuclei, ordered
pointwise. For $j\in\Nuc(H)$ write $\Fix(j)=\{a:ja=a\}$; it is closed under
meets and under $b\to(-)$ for every $b\in H$, i.e.\ $b\to a\in\Fix(j)$ whenever
$a\in\Fix(j)$, since $j(b\to a)\le jb\to ja=jb\to a\le b\to a$. Conversely a
subset $S\subseteq H$ closed under finite meets and under $b\to(-)$ is the
fixed set of at most one nucleus, $j_S(a)=\bigwedge\{s\in S:a\le s\}$ when this
meet exists for all $a$ (always if $H$ is complete or finite).

\begin{lem}\label{lem:join}
Let $j_1,j_2\in\Nuc(H)$. If the meet
$ja:=\bigwedge\{s\in\Fix(j_1)\cap\Fix(j_2):a\le s\}$ exists for every $a\in H$,
then $j$ is a nucleus, $\Fix(j)=\Fix(j_1)\cap\Fix(j_2)$, and $j$ is the least
nucleus above $j_1$ and $j_2$ in $\Nuc(H)$. If $H$ is complete the meet always
exists, and $j=j_1\vee j_2$ in the complete lattice $\Nuc(H)$.
\end{lem}
\begin{proof}
$S=\Fix(j_1)\cap\Fix(j_2)$ is closed under meets and under $b\to(-)$, so $j=j_S$
is a nucleus with fixed set $S$ (the standard argument: $j$ is inflationary,
monotone and idempotent; $j(a\wedge b)\ge ja\wedge jb$ because
$a\to j(a\wedge b)\in S$ contains $b$, hence $jb$, so $a\wedge jb\le j(a\wedge b)$,
then $jb\to j(a\wedge b)\in S$ contains $a$, hence $ja$). Any nucleus $k\ge j_1,j_2$
has $\Fix(k)\subseteq S$, so $ja\le ka$; and $j\ge j_i$ since $\Fix(j)\subseteq\Fix(j_i)$.
For complete $H$ see \cite[II.2.3]{Johnstone}.
\end{proof}

\begin{lem}\label{lem:lfp}
In a complete Heyting algebra, $(j_1\vee j_2)(a)$ is the least fixed point of
$X\mapsto a\vee j_1X\vee j_2X$, and it is $\bigvee_\alpha(j_1j_2)^\alpha(a)$ over
the transfinite iterates (Macnab \cite{Macnab}); on a finite algebra it is
$(j_1j_2)^n(a)$ for some $n$.
\end{lem}
\begin{proof}
A fixed point $X$ of the operator is above $a$ and fixed by $j_1$ and $j_2$,
so $X\ge(j_1\vee j_2)(a)$; and $(j_1\vee j_2)(a)$ is a fixed point. The iterates
$(j_1j_2)^\alpha(a)$ form an increasing chain of elements below
$(j_1\vee j_2)(a)$ which stabilizes at a common fixed point.
\end{proof}

\section{The logic $\LLJ$}\label{sec:logic}

Formulas are built from propositional variables with $\wedge,\vee,\to,\bot,\top$
and three unary connectives $\laxa,\laxb,\lax$. For a unary connective $\square$
the \emph{lax axioms} are
\[
\varphi\to\square\varphi,\qquad\square\square\varphi\to\square\varphi,\qquad
\square\varphi\wedge\square\psi\to\square(\varphi\wedge\psi),
\]
with the rule: from $\varphi\to\psi$ infer $\square\varphi\to\square\psi$.
$\PLL$ is intuitionistic propositional logic with one lax connective
\cite{FM}, $\PLLt$ has two, $\laxa$ and $\laxb$.

\begin{defi}\label{def:LJ}
$\LLJ$ is intuitionistic propositional logic with the lax axioms and rules for
$\laxa$, $\laxb$ and $\lax$, the axioms $\laxa\varphi\to\lax\varphi$ and
$\laxb\varphi\to\lax\varphi$, and the induction scheme (JL).
\end{defi}

An \emph{$\LLJ$-algebra} is a Heyting algebra $H$ with nuclei $j_1,j_2,j$ such
that $j\ge j_1$, $j\ge j_2$ and, for all $a,b\in H$,
\[
(j_1b\to b)\wedge(j_2b\to b)\wedge(a\to b)\le ja\to b.\tag{JL$'$}
\]
Since (JL$'$) is an identity, $\LLJ$-algebras form a variety. $\LLJ$ is
algebraizable in the sense of Blok and Pigozzi \cite{BP} with this variety as
equivalent algebraic semantics, with defining equation $\varphi\approx\top$
and equivalence formulas $\{\varphi\to\psi,\ \psi\to\varphi\}$: this is the
algebraizing pair of intuitionistic propositional logic, and it remains one for
an axiomatic extension by unary connectives $\square$ that are
\emph{congruential}, i.e.\ for which $\varphi\leftrightarrow\psi$ yields
$\square\varphi\leftrightarrow\square\psi$; the monotonicity rule gives this
for $\laxa,\laxb,\lax$. The axioms translate into the identities defining
$\LLJ$-algebras, (JL) into (JL$'$).

\begin{prop}\label{prop:least}
In an $\LLJ$-algebra, $ja$ is the least element of $\Fix(j_1)\cap\Fix(j_2)$
above $a$; hence $j$ is the least nucleus above $j_1$ and $j_2$, and
$\Fix(j)=\Fix(j_1)\cap\Fix(j_2)$. Conversely, if $H$ is a Heyting algebra with
nuclei $j_1,j_2$ for which the meets of Lemma~\ref{lem:join} exist (in
particular if $H$ is complete), then $(H,j_1,j_2,j_1\vee j_2)$ is an
$\LLJ$-algebra.
\end{prop}
\begin{proof}
Let $b\ge a$ be fixed by $j_1$ and $j_2$. The left side of (JL$'$) is $\top$, so
$ja\le b$. As $j\ge j_1,j_2$, $\Fix(j)\subseteq\Fix(j_1)\cap\Fix(j_2)$, and
$ja$ is fixed by $j_1,j_2$; so $ja$ is the least common fixed point above $a$
and the rest follows as in Lemma~\ref{lem:join}. For the converse, let
$j=j_1\vee j_2$ and $c\le(j_1b\to b)\wedge(j_2b\to b)\wedge(a\to b)$. Then
$c\wedge j_ib\le b$, so $j_ib\le c\to b$ and
$j_i(c\to b)\le j_ic\to j_ib\le c\to j_ib\le c\to(c\to b)=c\to b$: the element
$c\to b$ is fixed by $j_1$ and $j_2$, and it is above $a$ because $c\le a\to b$.
Hence $ja\le c\to b$, i.e.\ $c\le ja\to b$.
\end{proof}

\begin{thm}\label{thm:complete}
$\LLJ\vdash\varphi$ if and only if $\varphi$ is valid in every $\LLJ$-algebra.
Every complete Heyting algebra with two nuclei $j_1,j_2$, with $\lax$
interpreted as $j_1\vee j_2$, is an $\LLJ$-algebra; in particular every
Grothendieck topos equipped with two Lawvere--Tierney topologies and their join
validates $\LLJ$ on its subterminal objects, where the topologies act as nuclei
and the join of topologies as the join of nuclei.
\end{thm}
\begin{proof}
Soundness for $\LLJ$-algebras is Proposition~\ref{prop:least} together with
the lax axioms, and completeness is the Lindenbaum--Tarski algebra, which is an
$\LLJ$-algebra because the axioms are identities in it. Complete Heyting
algebras with two nuclei are $\LLJ$-algebras by Proposition~\ref{prop:least}.
For the topos statement, the subterminal objects of a Grothendieck topos form a
complete Heyting algebra, a Lawvere--Tierney topology restricts to a nucleus on
it, and the join of two topologies restricts to the join of the two nuclei,
since the restriction $\Omega\to\mathrm{Sub}(1)$ preserves the lattice
structure of local operators (both are computed by the same transfinite
iteration of the composite, Lemma~\ref{lem:lfp}).
\end{proof}

\begin{rem}\label{rem:frames}
Theorem~\ref{thm:complete} leaves open whether $\LLJ$ is complete for the
\emph{complete} Heyting algebras with two nuclei, i.e.\ whether every
$\LLJ$-algebra embeds, as a Heyting algebra with nuclei $j_1,j_2$, into a
complete one in which the join of the extended nuclei restricts to $j$.
Nuclei extend to the MacNeille completion (Goldblatt \cite[\S4]{Goldblatt}) and
to the canonical extension, but the join of the extensions may be larger than
the extension of the join, since the completion has more common fixed points
below $ja$ than the original algebra has. The finite model property would
settle this, as finite algebras are complete (Question~\ref{que:fmp}); the
poset semantics of the next section is then the intended one.
\end{rem}

\section{Finite posets and finite sites}\label{sec:poset}

Let $P$ be a finite poset and $H=\Up(P)$ its Heyting algebra of up-sets. For
$C\subseteq P$ put
\[
\lax_C U=\{x\in P:\ \uparrow\!x\cap C\subseteq U\},
\]
``$U$ holds at every point of $C$ above $x$''.

\begin{thm}\label{thm:poset}
\begin{enumerate}
\item[(i)] $\lax_C$ is a nucleus on $\Up(P)$, $C\mapsto\lax_C$ is an
order-reversing bijection from the subsets of $P$ onto $\Nuc(\Up(P))$, and
$\lax_C\wedge\lax_{C'}=\lax_{C\cup C'}$, $\lax_C\vee\lax_{C'}=\lax_{C\cap C'}$.
\item[(ii)] $\lax_C\vee\lax_{C'}=(\lax_C\lax_{C'})^N$ for every $N$ with $2N>|P|$,
in particular for $N=|P|$; and for every $N$ there is a poset with $2N$ points
on which $N-1$ iterations do not suffice.
\item[(iii)] For $2N>|P|$ the scheme $(\laxa\laxb)^{N+1}\varphi\to(\laxa\laxb)^N\varphi$
is valid on $(\Up(P),\lax_C,\lax_{C'})$, and $\LLJ$ plus this scheme proves
$\lax\varphi\leftrightarrow(\laxa\laxb)^N\varphi$.
\end{enumerate}
\end{thm}
\begin{proof}
(i) $\lax_C$ preserves intersections and is inflationary and idempotent
(if $x\in\lax_C\lax_CU$ and $c\in C$ is above $x$, then $c\in\lax_CU$ and $c\in C$
above $c$ gives $c\in U$). For the bijection we use the dictionary of
\cite[\S2]{Finitejoin}: the nuclei on $\Up(P)$ are the Lawvere--Tierney
topologies of the presheaf topos on $P$, which correspond to Grothendieck
topologies on $P$, hence (Kelly--Lawvere, Marqu\`es, \cite[Thm.~5.1]{Postnikov})
to idempotent two-sided ideals of the category $P$, and for a finite poset an
idempotent ideal is generated by the identities it contains, i.e.\ by a subset
$C\subseteq P$: the ideal $D_C=\{(x,y):x\le c\le y\text{ for some }c\in C\}$.
The nucleus of $D_C$ on sieves (down-sets) is $S\mapsto\{y:D_C(-,y)\subseteq S\}
=\{y:\downarrow\!y\cap C\subseteq S\}$, which is $\lax_C$ after passing from
down-sets to up-sets. Meets of nuclei are pointwise, and
$\lax_CU\cap\lax_{C'}U=\lax_{C\cup C'}U$. Joins: $\Fix(\lax_C)=\{U:\forall x\,
(\uparrow\!x\cap C\subseteq U\Rightarrow x\in U)\}$, and
$\Fix(\lax_C)\cap\Fix(\lax_{C'})=\Fix(\lax_{C\cap C'})$: the inclusion
$\supseteq$ holds because $\uparrow\!x\cap C\subseteq U$ implies
$\uparrow\!x\cap(C\cap C')\subseteq U$; for $\subseteq$, let $U$ be fixed by both
and $x\notin U$ with $\uparrow\!x\cap C\cap C'\subseteq U$; choose $x$ maximal
with these properties; as $U$ is $\lax_C$-fixed there is $c\in C$ above $x$
with $c\notin U$, and $c\notin C'$ (else $c\in U$); as $c\notin U$ is
$\lax_{C'}$-fixed-violating, there is $c'\in C'$ above $c$ with $c'\notin U$,
$c'\notin C$; then $c'>x$ has the same properties, contradicting maximality
(note $\uparrow\!c'\cap C\cap C'\subseteq\uparrow\!x\cap C\cap C'\subseteq U$).
By Lemma~\ref{lem:join} the join is $\lax_{C\cap C'}$.

(ii) By induction, $x\in(\lax_C\lax_{C'})^nU$ iff every alternating chain
$x\le c_1\le c_1'\le c_2\le\dots\le c_n'$ with $c_i\in C$, $c_i'\in C'$ ends in
$U$. If $x\notin\lax_{C\cap C'}U$, there is $z\in C\cap C'$ above $x$ with
$z\notin U$, and the constant chain $c_i=c_i'=z$ shows $x\notin(\lax_C\lax_{C'})^nU$
for all $n$. Conversely let $x\in\lax_{C\cap C'}U$ and $2N>|P|$, and let
$x\le c_1\le c_1'\le\dots\le c_N'$ be an alternating chain. Its $2N>|P|$
entries cannot all be distinct, and since the chain is monotone two adjacent
entries coincide: $c_i=c_i'$ or $c_i'=c_{i+1}$, a point of $C\cap C'$ above
$x$, hence in $U$; the chain stays above it and ends in $U$. For the second
claim take $P$ the chain $1<2<\dots<2N$, $C$ the odd points together with
$2N$, $C'$ the even points, and $U=\{2N\}$: then $C\cap C'=\{2N\}$ and
$1\in\lax_{C\cap C'}U$, while the chain $1\le1\le2\le3\le4\le\dots\le2N-2$
of $2(N-1)$ entries ends outside $U$, so $1\notin(\lax_C\lax_{C'})^{N-1}U$.

(iii) The first claim is (ii). In $\LLJ$ plus the scheme, $\theta:=(\laxa\laxb)^N\varphi$
satisfies $\varphi\to\theta$, $\laxa\theta\to\theta$ (as $\laxa\theta=\laxa(\laxa\laxb)^N\varphi$
is provably equivalent to $(\laxa\laxb)^N\varphi$ by idempotence of $\laxa$) and
$\laxb\theta\to\theta$ (as $\laxb(\laxa\laxb)^N\varphi\to\laxa\laxb(\laxa\laxb)^N\varphi=(\laxa\laxb)^{N+1}\varphi\to\theta$
by the scheme), so (JL) gives $\lax\varphi\to\theta$; and $\theta\to\lax\varphi$
since $\laxa,\laxb\le\lax$ and $\lax$ is idempotent.
\end{proof}

\begin{rem}
For a finite category $S$ in place of a poset, the nuclei on the subterminal
objects of $\Presh(S)$ are again the Grothendieck topologies, i.e.\ the
idempotent ideals $D$, with $\lax_D$ given by $D(-,c)\subseteq S$; the join
of $\lax_D$ and $\lax_E$ is $\lax_{I'}$ for the largest idempotent ideal
$I'\subseteq D\cap E$, and it is reached by $2m$ alternating composites, where
$m$ is the length of the chain $DE\supseteq DEDE\supseteq\cdots$
\cite[\S2]{Finitejoin}. The sameness-lattice theorem of \cite{Finitejoin},
$W_D\vee W_E=W_{I'}$, is the statement that the maps of presheaves
$f$ with $\lax_{I'}\llbracket f\text{ is an isomorphism}\rrbracket=\top$ are
exactly the two-out-of-three combinations of maps $g$ with
$\lax_D\llbracket g\text{ is an isomorphism}\rrbracket=\top$ or
$\lax_E\llbracket g\text{ is an isomorphism}\rrbracket=\top$; it is a
statement about maps and does not reduce to the propositional logic $\LLJ$.
\end{rem}

\section{Inexpressibility}\label{sec:inexp}

\begin{thm}\label{thm:inexp}
There is no formula $\chi$ of $\PLLt$ (in the connectives $\wedge,\vee,\to,
\bot,\top,\laxa,\laxb$ and propositional variables) such that
$\lax\bot\leftrightarrow\chi$ is valid in all $\LLJ$-algebras, nor in all
complete chains with two nuclei. Consequently no $\chi(p)$ defines $\lax p$,
and for every $n$ the formula $\lax\bot\leftrightarrow(\laxa\laxb)^n\bot$ fails
in some complete chain.
\end{thm}
\begin{proof}
Let $H$ be the chain $\omega+2=\{0<1<2<\dots<\omega<\omega+1\}$, a complete
Heyting algebra with $a\to b=\top$ if $a\le b$ and $a\to b=b$ otherwise. On a
chain every inflationary idempotent monotone map is a nucleus. Let
\[
j_1(2k)=2k+1,\quad j_1(2k+1)=2k+1,\qquad j_2(2k+1)=2k+2,\quad j_2(2k)=2k,
\]
and $j_1(\omega)=j_2(\omega)=\omega$, $j_1(\omega+1)=j_2(\omega+1)=\omega+1$.
Then $\Fix(j_1)=\{\text{odd}\}\cup\{\omega,\omega+1\}$,
$\Fix(j_2)=\{\text{even}\}\cup\{\omega,\omega+1\}$, so by Lemma~\ref{lem:join}
$(j_1\vee j_2)(a)=\omega$ for all $a\le\omega$, in particular
$(j_1\vee j_2)(\bot)=\omega$. The set $A=\mathbb N\cup\{\omega+1\}$ contains
$\bot=0$ and $\top=\omega+1$, is closed under $j_1,j_2$, under $\wedge,\vee$,
and under $\to$ (as $a\to b\in\{\top,b\}$). If $\chi$ had the stated property,
then, whatever values in $H$ are assigned to its variables --- assign $\bot$ to
all of them --- the value of $\chi$ would be $\omega$; but the value of a
formula all of whose variables take values in $A$ lies in $A$, and
$\omega\notin A$. The second statement follows by substituting $\bot$ for $p$;
the third because $(j_1j_2)^n(0)=2n\ne\omega$.
\end{proof}

\begin{rem}
The join is implicitly definable: in any Heyting algebra with nuclei $j_1,j_2$
there is at most one nucleus $j$ satisfying the $\LLJ$-axioms
(Proposition~\ref{prop:least}). Theorem~\ref{thm:inexp} is thus a failure of
explicit definability for a connective, of the same kind as the
inexpressibility of common knowledge by finite iterations of ``everybody
knows'' \cite{HM}: there $C\varphi=\bigwedge_nE^n\varphi$ is a greatest fixed
point of a deflationary operator and the induction axiom
$C(\varphi\to E\varphi)\to(\varphi\to C\varphi)$ expresses it; here
$\lax\varphi=\bigvee_n(\laxa\laxb)^n\varphi$ (on complete algebras, with
transfinite iteration in general) is a least fixed point of an inflationary
operator and (JL) expresses it. In the chain of Theorem~\ref{thm:inexp} the
join needs exactly $\omega$ iterations.
\end{rem}

\begin{cor}[uniqueness without definability]\label{cor:beth}
Let $\LLJ^{\lax,\lax'}$ be $\LLJ$ with a second copy $\lax'$ of the join
connective, satisfying the same axioms relative to $\laxa,\laxb$. Then
$\LLJ^{\lax,\lax'}\vdash\lax\varphi\leftrightarrow\lax'\varphi$ for all
$\varphi$, while no formula $\chi(p)$ of $\PLLt$ satisfies
$\LLJ\vdash\lax p\leftrightarrow\chi(p)$. Thus the join connective is
implicitly but not explicitly definable over $\PLLt$: the Beth definability
property for connectives fails for this extension.
\end{cor}
\begin{proof}
In an algebra of the extended logic both $j$ and $j'$ send $a$ to the least
common fixed point of $j_1,j_2$ above $a$ (Proposition~\ref{prop:least}), so
$j=j'$ and the equivalence is valid, hence provable by algebraic completeness.
The second claim is Theorem~\ref{thm:inexp}.
\end{proof}

\begin{rem}
Everything extends to $r$ lax modalities $\lax_1,\dots,\lax_r$ and their
join $\lax$: the axioms are $\lax_i\varphi\to\lax\varphi$ and
$(\lax_1\psi\to\psi)\to\dots\to(\lax_r\psi\to\psi)\to(\varphi\to\psi)\to(\lax\varphi\to\psi)$,
the algebras are those in which $\lax a$ is the least common fixed point of
$\lax_1,\dots,\lax_r$ above $a$, on a finite poset the join of
$\lax_{C_1},\dots,\lax_{C_r}$ is $\lax_{C_1\cap\dots\cap C_r}$, and the
chain of Theorem~\ref{thm:inexp} with $r-2$ identity nuclei added shows
inexpressibility.
\end{rem}

\begin{que}\label{que:fmp}
Does $\LLJ$ have the finite model property, i.e.\ is every non-theorem
refuted on a finite poset with two subsets $C,C'$ (Theorem~\ref{thm:poset})?
Is $\LLJ$ decidable, and is it complete for complete Heyting algebras
(Remark~\ref{rem:frames})? For the logic of common knowledge the first two
hold, by finite canonical models \cite{HM}. A natural attempt is an algebraic
filtration: the values of the subformulas of a non-theorem generate a finite
sublattice $D$ of the refuting algebra, closed under $d\mapsto d\to s$ for
$s$ in the finite set $S$ of values of modal subformulas; on $D$ the closure
operators whose fixed sets are the meet-closures of $\{d\to s:d\in D,\ s\in S_i\}$
are nuclei agreeing with $j_i$ on subformula values, and their join agrees with
$j$ on subformula values. The gap is the finiteness of $D$: for one $s$ the closure of a finite set
under $\wedge$, $\vee$ and $d\mapsto d\to s$ is finite (it lives in the
Boolean algebra of elements of the form $x\to s$, as for pseudocomplemented
lattices), but for two elements $s,s'$ we do not know whether it is, and in
the free Heyting algebra on $p,s,s'$ the number of pairwise inequivalent terms
produced in three rounds of closure already exceeds $30{,}000$, so the
argument does not apply as it stands.
\end{que}

\begin{figure}[ht]
\[
\begin{tikzcd}[column sep=large]
\varphi\arrow[r,"\laxa"]\arrow[rrr,bend left=25,"\lax=\laxa\vee\laxb"] & \laxa\varphi\arrow[r,"\laxb"] & \laxb\laxa\varphi\arrow[r,"\cdots"] & (\laxa\laxb)^N\varphi
\end{tikzcd}
\]
\caption{The join as the limit of alternating iterates: finite on a finite
poset ($N=|P|$, Theorem~\ref{thm:poset}), transfinite in general, and not a
$\PLLt$-formula (Theorem~\ref{thm:inexp}).}
\end{figure}

\section*{Funding}
No funding was received for this work.

\section*{Declarations}
\noindent\textbf{CRediT authorship contribution statement.}\quad Ryota Shibusawa:
Conceptualization, Methodology, Formal analysis, Investigation, Software,
Validation, Visualization, Writing -- original draft, Writing -- review \& editing.

\noindent\textbf{Competing interests.}\quad The author declares that he has no
competing interests.

\begin{thebibliography}{99}
\bibitem{BBCGGJ} G.~Bezhanishvili, N.~Bezhanishvili, L.~Carai, D.~Gabelaia,
S.~Ghilardi, M.~Jibladze, \emph{Diego's theorem for nuclear implicative
semilattices}, Indag. Math. \textbf{32} (2021), 498--535.
\bibitem{BH} G.~Bezhanishvili, W.~H.~Holliday, \emph{Locales, nuclei, and
Dragalin frames}, in: Advances in Modal Logic 11, College Publications, 2016,
177--196.
\bibitem{BP} W.~J.~Blok, D.~Pigozzi, \emph{Algebraizable logics}, Mem. Amer.
Math. Soc. \textbf{396} (1989).
\bibitem{Escardo} M.~H.~Escard\'o, \emph{Joins in the frame of nuclei}, Appl.
Categ. Structures \textbf{11} (2003), 117--124.
\bibitem{FM} M.~Fairtlough, M.~Mendler, \emph{Propositional lax logic}, Inform.
and Comput. \textbf{137} (1997), 1--33.
\bibitem{Goldblatt} R.~Goldblatt, \emph{Grothendieck topology as geometric
modality}, Z. Math. Logik Grundlag. Math. \textbf{27} (1981), 495--529.
\bibitem{GoldblattCover} R.~Goldblatt, \emph{Cover semantics for quantified lax
logic}, J. Logic Comput. \textbf{21} (2011), 1035--1063.
\bibitem{HM} J.~Y.~Halpern, Y.~Moses, \emph{Knowledge and common knowledge in a
distributed environment}, J. ACM \textbf{37} (1990), 549--587.
\bibitem{Johnstone} P.~T.~Johnstone, \emph{Stone Spaces}, Cambridge University
Press, 1982.
\bibitem{KL89} G.~M.~Kelly, F.~W.~Lawvere, \emph{On the complete lattice of
essential localizations}, Bull. Soc. Math. Belg. S\'er. A \textbf{41} (1989),
289--319.
\bibitem{Macnab} D.~S.~Macnab, \emph{Modal operators on Heyting algebras},
Algebra Universalis \textbf{12} (1981), 5--29.
\bibitem{Marques25} J.~Marqu\`es, \emph{A criterion for categories on which
every Grothendieck topology is rigid}, Appl. Categ. Structures \textbf{33}
(2025), art.\ 37.
\bibitem{Postnikov} R.~Shibusawa, \emph{Postnikov truncations as joins of sheaf
and homotopy equivalences}, preprint (2026), archived at DOI 10.5281/zenodo.22956181.
\bibitem{Finitejoin} R.~Shibusawa, \emph{Grothendieck topologies on a finite
category form a sublattice of the lattice of two-out-of-three classes},
preprint (2026), archived at DOI 10.5281/zenodo.23204943.
\end{thebibliography}

\end{document}
